% =====================================================================
%  Soutenance du projet final Data Engineer, Spark et Scala
%  Groupe 9 : ABOUTA Eudoxie, BAMBA Issouf, AHOGA Josias
%
%  Compilation : pdfLaTeX (Overleaf, bouton Recompile)
%  Aucune image externe, aucun paquet exotique.
% =====================================================================
\documentclass[aspectratio=169,11pt]{beamer}

\usepackage[T1]{fontenc}
\usepackage[utf8]{inputenc}
\usepackage[french]{babel}
\usepackage{cmbright}
\usepackage{booktabs}
\usepackage{array}
\usepackage{xcolor}
\usepackage{tikz}
\usepackage{listings}

% ---------------------------------------------------------------- couleurs
\definecolor{indigo}{HTML}{2D288C}
\definecolor{indigofonce}{HTML}{221E69}
\definecolor{grisbloc}{HTML}{EFEFF5}
\definecolor{encre}{HTML}{1A1C2B}
\definecolor{gris}{HTML}{6A6F82}
\definecolor{vert}{HTML}{1E8E5A}
\definecolor{vertpale}{HTML}{D6EDE1}
\definecolor{rouge}{HTML}{B5322E}
\definecolor{ambre}{HTML}{C98A1B}

% ---------------------------------------------------------------- theme
\usetheme{default}
\usecolortheme{default}
\setbeamertemplate{navigation symbols}{}
\setbeamercolor{structure}{fg=indigo}
\setbeamercolor{normal text}{fg=encre}
\setbeamersize{text margin left=9mm,text margin right=9mm}

\setbeamertemplate{itemize item}{\color{indigo}\rule[0.35ex]{1.7mm}{0.8pt}}
\setbeamertemplate{itemize subitem}{\color{gris}\rule[0.35ex]{1.2mm}{0.6pt}}
\setbeamertemplate{section in toc}{\leavevmode\color{indigo}\rule[0.35ex]{1.7mm}{0.8pt}\ \color{encre}\inserttocsection\par}
\setbeamerfont{section in toc}{size=\normalsize}

% Bandeau de titre pleine largeur
\setbeamercolor{frametitle}{bg=indigo,fg=white}
\setbeamerfont{frametitle}{size=\large}
\setbeamertemplate{frametitle}{%
  \nointerlineskip
  \begin{beamercolorbox}[wd=\paperwidth,ht=3.2ex,dp=1.6ex,leftskip=9mm,rightskip=9mm]{frametitle}
    \usebeamerfont{frametitle}\insertframetitle
  \end{beamercolorbox}%
  \vskip3mm}

% Blocs a en tete plein
\setbeamercolor{block title}{bg=indigo,fg=white}
\setbeamercolor{block body}{bg=grisbloc,fg=encre}
\setbeamerfont{block title}{size=\normalsize,series=\mdseries}
\setbeamertemplate{blocks}[default]

% Pied de page a quatre cellules
\setbeamercolor{cellule}{bg=indigo,fg=white}
\setbeamercolor{cellulesombre}{bg=indigofonce,fg=white}
\newcommand{\pieddepage}{%
  \setbeamertemplate{footline}{%
    \hbox{%
      \begin{beamercolorbox}[wd=.26\paperwidth,ht=2.5ex,dp=1.2ex,leftskip=9mm]{cellulesombre}
        \scriptsize GROUPE 9
      \end{beamercolorbox}%
      \begin{beamercolorbox}[wd=.44\paperwidth,ht=2.5ex,dp=1.2ex,center]{cellule}
        \scriptsize EcommerceAnalytics
      \end{beamercolorbox}%
      \begin{beamercolorbox}[wd=.20\paperwidth,ht=2.5ex,dp=1.2ex,center]{cellulesombre}
        \scriptsize 2025-2026
      \end{beamercolorbox}%
      \begin{beamercolorbox}[wd=.10\paperwidth,ht=2.5ex,dp=1.2ex,center]{cellule}
        \scriptsize \insertframenumber/\inserttotalframenumber
      \end{beamercolorbox}}%
    \vskip0pt}}
\newcommand{\sanspied}{\setbeamertemplate{footline}{}}
\pieddepage

% ---------------------------------------------------------------- code
\lstdefinestyle{scala}{
  basicstyle=\ttfamily\scriptsize\color{encre},
  backgroundcolor=\color{grisbloc},
  keywordstyle=\color{indigo}\bfseries,
  stringstyle=\color{vert},
  commentstyle=\color{gris}\itshape,
  showstringspaces=false,
  breaklines=true,
  frame=none,
  xleftmargin=4mm, xrightmargin=2mm,
  framexleftmargin=4mm, framexrightmargin=2mm,
  aboveskip=2mm, belowskip=2mm,
  columns=fullflexible,
  keepspaces=true
}
\lstset{style=scala}

% ---------------------------------------------------------------- outils
% Legende grise sous un element
\newcommand{\legende}[1]{{\color{gris}\footnotesize #1}}

% Barre horizontale : \barre{fraction}{couleur}{etiquette}
% Largeur en centimetres reglable par \renewcommand{\lbarre}{...}
\newcommand{\lbarre}{5}
\newcommand{\barre}[3]{%
  \begin{tikzpicture}[baseline=-0.7ex]
    \fill[grisbloc] (0,0) rectangle (\lbarre,0.26);
    \fill[#2] (0,0) rectangle ({#1*\lbarre},0.26);
  \end{tikzpicture}\;{\scriptsize #3}}

% ---------------------------------------------------------------- meta
\title{Analyse de données e commerce}
\author{Groupe 9}
\date{2025-2026}

\begin{document}

% ================================================================ 1
\sanspied
\begin{frame}[plain]
  \vspace{3mm}
  \begin{center}
  \begin{tikzpicture}
    \node[fill=indigo,text=white,rounded corners=6pt,inner sep=5mm,
          text width=\dimexpr\textwidth-12mm\relax,align=center]
      {\Large Analyse de données e commerce\\[3mm]
       \normalsize Un pipeline Spark qui explique ses rejets, mesure la fidélité\\
       et publie lui même son tableau de bord};
  \end{tikzpicture}

  \vspace{14mm}
  Projet final Data Engineer, Spark et Scala\\
  Dakar Institute of Technology

  \vspace{6mm}
  {\bfseries ABOUTA EUDOXIE \quad BAMBA ISSOUF \quad AHOGA JOSIAS}

  \vspace{6mm}
  Année académique 2025-2026
  \end{center}

  \note{Collectif, 30 secondes. Se présenter, annoncer le plan. Préciser
  d'emblée que les données sont synthétiques et que la démonstration se
  termine sur un tableau de bord généré par le pipeline lui même.}
\end{frame}
\pieddepage

% ================================================================ 2
\begin{frame}{Plan}
  \vspace{3mm}
  \begin{center}
  \begin{minipage}{0.80\textwidth}
    \setlength{\parskip}{2.6mm}
    \tableofcontents
  \end{minipage}
  \end{center}

  \note{Collectif, 20 secondes. Annoncer les trois exposés individuels et la
  démonstration finale. Ne pas commenter chaque section.}
\end{frame}

\section{Contexte et problématique}

% ================================================================ 3
\begin{frame}{Contexte et problématique}
  \begin{itemize}\itemsep3mm
    \item Quatre fichiers de formats différents, volontairement sales, décrivent
          la même activité commerciale.
    \item 138\,047 transactions tiennent sur un poste de travail : la difficulté
          n'est pas la taille, c'est la fiabilité de ce qu'on en tire.
    \item Un montant invalide fausse le chiffre d'affaires, la commission du
          marchand et son classement.
  \end{itemize}

  \vspace{4mm}
  \begin{block}{Traçabilité}
    Chaque ligne écartée porte son motif. Tout chiffre présenté se remonte
    jusqu'à sa ligne source.
  \end{block}

  \note{Collectif, 45 secondes. La qualité vient avant le chiffre d'affaires.
  Annoncer que l'égalité de réconciliation sera montrée en détail par Eudoxie.}
\end{frame}

% ================================================================ 4
\begin{frame}{Les quatre sources}
  \vspace{1mm}
  \begin{center}
  \small
  \begin{tabular}{@{}l r l p{5.6cm}@{}}
    \toprule
    \textbf{Fichier} & \textbf{Lignes} & \textbf{Format} & \textbf{Stratégie de lecture} \\
    \midrule
    \texttt{transactions.csv} & 138\,047 & CSV     & Schéma explicite, mode \texttt{FAILFAST} \\
    \texttt{users.json}       &  12\,000 & JSON    & Schéma typé, champ tableau imbriqué \\
    \texttt{products.parquet} &   6\,000 & Parquet & Répertoire de douze fichiers \\
    \texttt{merchants.csv}    &     600  & CSV     & Inférence, puis conversions explicites \\
    \bottomrule
  \end{tabular}
  \end{center}

  \vspace{5mm}
  \begin{block}{Configuration externalisée}
    Chemins, seuils et paramètres Spark proviennent de
    \texttt{application.conf}. Aucune valeur codée en dur.
  \end{block}

  \note{Collectif, 60 secondes. Chaque format impose une stratégie de lecture
  différente. Si le jury demande pourquoi les données restent hors du JAR :
  pour être remplaçables et accessibles aux exécuteurs sur un cluster.}
\end{frame}

\section{Architecture et organisation}

% ================================================================ 5
\begin{frame}{Architecture en médaillon}
  \vspace{1mm}
  \begin{columns}[t,onlytextwidth]
    \column{0.315\textwidth}
      \begin{block}{Bronze}
        \footnotesize
        Lecture multi format\\
        Schéma explicite et typage\\
        \texttt{FAILFAST} sur l'illisible\\[2mm]
        \legende{DataIngestion}
      \end{block}
    \column{0.315\textwidth}
      \begin{block}{Silver}
        \footnotesize
        Rejets conservés et motivés\\
        Jointures expliquées\\
        Fenêtres et UDF temporelle\\[2mm]
        \legende{DataValidation, TimeFeatures}
      \end{block}
    \column{0.315\textwidth}
      \begin{block}{Gold}
        \footnotesize
        Quatorze rapports métier\\
        Tableau de bord HTML\\
        Durées par étape\\[2mm]
        \legende{Analytics, DashboardReport}
      \end{block}
  \end{columns}

  \vspace{6mm}
  \begin{block}{Socle transverse}
    \texttt{application.conf}, cache et diffusion, \texttt{MainApp} pour
    l'orchestration et le chronométrage, conteneur et intégration continue.
  \end{block}

  \note{Collectif, 75 secondes. Le vocabulaire Bronze Silver Gold rend le
  lignage lisible sans expliquer Spark. Montrer MainApp : la session est créée
  une seule fois, et le finally libère le cache puis arrête Spark même en cas
  d'échec.}
\end{frame}

% ================================================================ 6
\begin{frame}{Organisation de l'équipe}
  \vspace{1mm}
  \begin{center}
  \small
  \begin{tabular}{@{}l l l l@{}}
    \toprule
    \textbf{Membre} & \textbf{Rôle} & \textbf{Parties} & \textbf{Relu par} \\
    \midrule
    ABOUTA Eudoxie & Ingestion et plateforme   & 1, 2 et 7 & Issouf et Josias \\
    BAMBA Issouf   & Transformations           & 3         & Eudoxie \\
    AHOGA Josias   & Analytique et performance & 4, 5 et 6 & Issouf et Eudoxie \\
    \bottomrule
  \end{tabular}
  \end{center}

  \vspace{5mm}
  \begin{columns}[t,onlytextwidth]
    \column{0.47\textwidth}
      \begin{block}{Séparation par fichier}
        \footnotesize Chaque fichier a un propriétaire unique.
      \end{block}
    \column{0.47\textwidth}
      \begin{block}{Relecture tracée}
        \footnotesize Chaque module est relu par un autre membre.
        Date, relecteur et remarques sont consignés.
      \end{block}
  \end{columns}

  \note{Collectif, 90 secondes. Expliquer les trois interfaces qui font tenir
  l'ensemble : Sources, ValidationResult et Enrichment. Ce sont elles qui
  permettent à chacun de travailler sans casser le module du voisin.}
\end{frame}

\section{Ingestion et qualité des données}

% ================================================================ 7
\begin{frame}[fragile]{Ingestion et typage}
  \legende{ABOUTA Eudoxie}

  \vspace{2mm}
\begin{lstlisting}
def transactions(): Dataset[Transaction] = read("transactions") {
  spark.read.option("header", true).option("mode", "FAILFAST")
    .schema(Encoders.product[Transaction].schema)
    .csv(path("transactions")).as[Transaction]
}
\end{lstlisting}

  \vspace{3mm}
  \begin{columns}[t,onlytextwidth]
    \column{0.47\textwidth}
      \begin{block}{Montant absent}
        \footnotesize Il reste nul jusqu'à la validation, il n'est jamais
        remplacé par zéro.
      \end{block}
    \column{0.47\textwidth}
      \begin{block}{Lecture immédiate}
        \footnotesize Une action \texttt{count} force la lecture dès
        l'ingestion.
      \end{block}
  \end{columns}

  \note{Eudoxie, 2 minutes 30. Montrer read, le try catch et le message d'erreur
  contextualisé. Expliquer pourquoi users.json est lu avec un schéma typé et non
  en inférence : le champ preferred categories est un tableau, l'inférence sur
  un fichier partiel donnerait un type instable.}
\end{frame}

% ================================================================ 8
\begin{frame}{Validation et rapport de qualité}
  \legende{ABOUTA Eudoxie}
  \renewcommand{\lbarre}{3.2}

  \vspace{2mm}
  \begin{columns}[t,onlytextwidth]
    \column{0.50\textwidth}
      \vspace{2mm}
      \barre{1.000}{vert}{Lues \quad 138\,047}\\[2.5mm]
      \barre{0.986}{indigo}{Valides \quad 136\,157}\\[2.5mm]
      \barre{0.904}{ambre}{Analysées \quad 124\,751}\\[4mm]
      \legende{Orphelines : 400 utilisateurs,\\300 produits, 250 marchands.}
    \column{0.46\textwidth}
      \small
      \begin{tabular}{@{}l r r@{}}
        \toprule
        \textbf{Jeu} & \textbf{Rejets} & \textbf{Taux} \\
        \midrule
        transactions & 1\,890 & 1,37 \% \\
        users        &   345  & 2,88 \% \\
        products     &   179  & 2,98 \% \\
        merchants    &    14  & 2,33 \% \\
        \bottomrule
      \end{tabular}
  \end{columns}

  \vspace{3mm}
  \begin{block}{Réconciliation des volumes}
    138\,047 = 1\,890 rejets + 11\,406 références non enrichissables + 124\,751 analysées.
  \end{block}

  \note{Eudoxie, 2 minutes 30. Ouvrir un CSV de rejets et commenter une ligne
  réelle. Les compteurs d'orphelins peuvent se recouper : une transaction peut
  pointer à la fois vers un utilisateur absent et un produit absent, elle est
  comptée dans les deux.}
\end{frame}

\section{Transformations et enrichissement}

% ================================================================ 9
\begin{frame}[fragile]{L'UDF extractTimeFeatures}
  \legende{BAMBA Issouf}

  \vspace{2mm}
\begin{lstlisting}
udf((value: String) => features(value, settings).orNull)

def dayPeriod(hour: Int, s: TimeSettings): String =
  if (hour < s.morningStart || hour >= s.nightStart) "Night"
  else if (hour < s.afternoonStart)                  "Morning"
  else if (hour < s.eveningStart)                    "Afternoon"
  else                                               "Evening"
\end{lstlisting}

  \vspace{3mm}
  \begin{columns}[t,onlytextwidth]
    \column{0.47\textwidth}
      \begin{block}{Robustesse}
        \footnotesize Une chaîne nulle, vide ou mal formée renvoie \texttt{null}.
        Le 29 février 2025 est rejeté.
      \end{block}
    \column{0.47\textwidth}
      \begin{block}{Découpage horaire}
        \footnotesize Morning 06h à 11h59, Afternoon 12h à 17h59, Evening 18h à
        21h59, Night 22h à 05h59.
      \end{block}
  \end{columns}

  \note{Issouf, 1 minute 30. Insister sur orNull : l'UDF retourne une structure
  ou null, elle ne lève jamais d'exception. Montrer le test qui vérifie les
  frontières à 21h59, 22h00 et minuit.}
\end{frame}

% ================================================================ 10
\begin{frame}[fragile]{Enrichissement et fenêtrage}
  \legende{BAMBA Issouf}

  \vspace{2mm}
\begin{lstlisting}
val glissante = Window.partitionBy("user_id").orderBy("epoch_seconds")
                      .rangeBetween(-7 * 86400L + 1L, 0L)

size(collect_set(col("transaction_date")).over(glissante))
\end{lstlisting}

  \vspace{3mm}
  \begin{columns}[t,onlytextwidth]
    \column{0.47\textwidth}
      \small
      \begin{tabular}{@{}l@{}}
        \texttt{transaction\_rank} \\
        \texttt{rolling\_amount\_7d} \\
        \texttt{active\_days\_7d} \\
        \texttt{days\_since\_previous} \\
        \texttt{age\_group} \\
      \end{tabular}\\[2mm]
      \legende{Colonnes ajoutées à chaque transaction.}
    \column{0.47\textwidth}
      \begin{block}{Fenêtre en secondes}
        \footnotesize La borne est en secondes, pas en nombre de lignes.
      \end{block}
  \end{columns}

  \note{Issouf, 2 minutes. Si le jury demande pourquoi les jointures sont en
  left : parce qu'une référence manquante doit être expliquée avant d'être
  exclue. La vérification require(retenues + rejetées == valides) garantit
  qu'aucune ligne n'est dupliquée ni perdue.}
\end{frame}

% ================================================================ 11
\begin{frame}{Détection de transactions suspectes}
  \legende{BAMBA Issouf}

  \vspace{3mm}
  \begin{columns}[t,onlytextwidth]
    \column{0.47\textwidth}
      \begin{itemize}\itemsep2.5mm
        \item Montant supérieur de 300 \% au panier moyen historique du client
        \item Achat entre 22h00 et 05h59
      \end{itemize}
    \column{0.47\textwidth}
      \begin{itemize}\itemsep2.5mm
        \item Moins de cinq minutes depuis l'achat précédent
        \item Moyen de paiement déclaré en \texttt{CRYPTO}
      \end{itemize}
  \end{columns}

  \vspace{6mm}
  \begin{block}{Chiffre clé}
    3\,760 transactions signalées, soit 3,01 \% du périmètre, avec au moins
    deux signaux sur quatre.
  \end{block}

  \note{Issouf, 1 minute 30. Point clé : la moyenne historique n'utilise que les
  achats strictement antérieurs, sans anticipation. Sans historique, le premier
  critère ne s'active pas, mais les trois autres peuvent suffire. La règle est
  pédagogique, ce n'est pas un modèle de fraude validé.}
\end{frame}

\section{Analytique métier et performance}

% ================================================================ 12
\begin{frame}{Performance par marchand}
  \legende{AHOGA Josias}
  \renewcommand{\lbarre}{6}

  \vspace{4mm}
  \barre{1.000}{vert}{TechStore 225 \quad 568\,073 EUR}\\[2.5mm]
  \barre{0.919}{indigo}{ElectroWorld 559 \quad 522\,206 EUR}\\[2.5mm]
  \barre{0.824}{ambre}{TechStore 112 \quad 468\,136 EUR}

  \vspace{6mm}
  \begin{columns}[t,onlytextwidth]
    \column{0.47\textwidth}
      \begin{block}{Classements}
        \footnotesize \texttt{dense\_rank} par catégorie et par région.
      \end{block}
    \column{0.47\textwidth}
      \begin{block}{Commission par ligne}
        \footnotesize Montant multiplié par le taux du marchand.
        TechStore 225 : 42\,889,54 EUR.
      \end{block}
  \end{columns}

  \note{Josias, 1 minute 30. Expliquer pourquoi le merchant id de la transaction
  fait foi pour le chiffre d'affaires, et non celui du catalogue produit :
  512 divergences sont signalées plutôt que corrigées en silence.}
\end{frame}

% ================================================================ 13
\begin{frame}{Cohortes et rétention}
  \legende{AHOGA Josias}
  \renewcommand{\lbarre}{3.4}

  \vspace{3mm}
  \begin{columns}[t,onlytextwidth]
    \column{0.50\textwidth}
      \barre{1.000}{vert}{M+0 \quad 100,0 \%}\\[2mm]
      \barre{0.613}{vert}{M+1 \quad \phantom{0}61,3 \%}\\[2mm]
      \barre{0.487}{ambre}{M+2 \quad \phantom{0}48,7 \%}\\[2mm]
      \barre{0.381}{ambre}{M+3 \quad \phantom{0}38,1 \%}\\[2mm]
      \barre{0.302}{rouge}{M+5 \quad \phantom{0}30,2 \%}
    \column{0.46\textwidth}
      \begin{block}{Horizons non observables}
        \footnotesize Une cohorte de novembre 2025 n'a pas de M+3 observable.
        Ces horizons sont exclus du classement.
      \end{block}
  \end{columns}

  \vspace{5mm}
  \begin{block}{Meilleure rétention à trois mois}
    Cohorte de juin 2025 : 118 clients encore actifs sur 310, soit 38,06 \%.
  \end{block}

  \note{Josias, 1 minute 30. Expliquer la grille : on génère explicitement les
  couples cohorte et mois observable, on joint l'activité en left, et un mois
  sans achat vaut zéro. Sans cette grille, un mois creux disparaîtrait au lieu
  de valoir zéro, et la courbe serait fausse.}
\end{frame}

% ================================================================ 14
\begin{frame}{Segmentation RFM}
  \legende{AHOGA Josias}

  \vspace{3mm}
  \begin{center}
  \small
  \begin{tabular}{@{}l r r r r@{}}
    \toprule
    \textbf{Segment calculé} & \textbf{VIP} & \textbf{Premium} & \textbf{Standard} & \textbf{Budget} \\
    \midrule
    Champions       & \textcolor{vert}{\textbf{451}} & 613 & 272 & 32 \\
    Clients fidèles & 266 & 854 & 1\,305 & 816 \\
    À risque        & 287 & 726 & 866 & 317 \\
    Nouveaux        &  18 & 146 & 618 & 724 \\
    Perdus          &   2 &  67 & 653 & \textcolor{rouge}{\textbf{1\,160}} \\
    \bottomrule
  \end{tabular}
  \end{center}

  \vspace{4mm}
  \begin{block}{Lecture de la matrice}
    451 Champions déclarés VIP, 1\,160 Perdus déclarés Budget.
    32 Budget se comportent en Champions, 2 VIP sont perdus.
  \end{block}

  \note{Josias, 1 minute 30. Scores par quintiles avec ntile(5), récence
  ordonnée à l'envers, égalités départagées par identifiant pour être
  reproductible. Assumer que Clients fidèles est une classe résiduelle.}
\end{frame}

% ================================================================ 15
\begin{frame}{Optimisations et mesure}
  \legende{AHOGA Josias}
  \renewcommand{\lbarre}{3}

  \vspace{3mm}
  \begin{columns}[t,onlytextwidth]
    \column{0.50\textwidth}
      \footnotesize
      \begin{itemize}\itemsep1.6mm
        \item \texttt{cache} sur les résultats réutilisés
        \item \texttt{persist} sur les transactions enrichies
        \item \texttt{broadcast} des trois référentiels
        \item exécution adaptative sur la session
        \item \texttt{unpersist} dans un \texttt{finally}
      \end{itemize}
    \column{0.46\textwidth}
      \barre{0.730}{vert}{\phantom{0}95,9 s}\\[2mm]
      \barre{0.822}{vert}{108,0 s}\\[2mm]
      \barre{0.936}{ambre}{122,9 s}\\[2mm]
      \barre{1.000}{rouge}{131,3 s}\\[3mm]
      \legende{Quatre exécutions du même code,\\même machine, même journée.}
  \end{columns}

  \vspace{4mm}
  \begin{block}{Variabilité mesurée}
    37 \% d'écart entre la plus rapide et la plus lente.
  \end{block}

  \note{Josias, 2 minutes. C'est la diapositive où l'on gagne des points en
  étant honnête. Un seul passage par mode ne sépare pas le cache de
  l'échauffement de la JVM. Une étape est même plus lente avec les
  optimisations, et nous conservons ce chiffre.}
\end{frame}

\section{Démonstration et industrialisation}

% ================================================================ 16
\begin{frame}[fragile]{Démonstration}
  \vspace{2mm}
\begin{lstlisting}
spark-submit --master 'local[2]' --driver-memory 3g
  --class com.ecommerce.analytics.MainApp
  dist/ecommerce-analytics.jar all
\end{lstlisting}

  \vspace{4mm}
  \begin{center}
  \small
  \begin{tabular}{@{}r l r l@{}}
    \toprule
    \textbf{124\,751} & transactions analysées & \textbf{10\,193} & clients uniques \\
    \textbf{49,90 M EUR} & de chiffre d'affaires & \textbf{399,98 EUR} & de panier moyen \\
    \textbf{3,01 \%} & de transactions signalées & \textbf{14} & rapports produits \\
    \bottomrule
  \end{tabular}
  \end{center}

  \vspace{3mm}
  \begin{block}{Restitution}
    Une commande produit les quatorze rapports et une page HTML autonome,
    en thème clair et en thème sombre.
  \end{block}

  \note{Collectif, 2 minutes. Lancer la commande, puis ouvrir
  output/dashboard.html par double clic. Commenter deux résultats métier : la
  concentration du chiffre d'affaires sur les marchands Electronics, et la
  cohérence entre segment RFM calculé et segment déclaré.}
\end{frame}

% ================================================================ 17
\begin{frame}{Reproductibilité}
  \vspace{2mm}
  \begin{columns}[t,onlytextwidth]
    \column{0.47\textwidth}
      \begin{block}{Environnement figé}
        \footnotesize \texttt{Dockerfile} et \texttt{docker-compose} figent
        Spark 3.5.6, Scala 2.12 et Java 17. \texttt{Makefile} sous Linux et
        macOS, \texttt{make.ps1} sous Windows.
      \end{block}
    \column{0.47\textwidth}
      \begin{block}{Vérification automatique}
        \footnotesize GitHub Actions compile, teste et empaquette à chaque
        \texttt{push}. Vingt cas de régression couvrent les frontières et les
        règles métier.
      \end{block}
  \end{columns}

  \vspace{6mm}
  \begin{center}
  \small
  \begin{tabular}{@{}r l @{\hspace{12mm}} r l @{\hspace{12mm}} r l@{}}
    \textbf{20 / 20} & tests au vert &
    \textbf{108 s} & pipeline complet &
    \textbf{0} & chemin codé en dur \\
  \end{tabular}
  \end{center}

  \note{Collectif, 90 secondes. Préciser honnêtement : la voie Docker est
  documentée à partir des images officielles, elle n'a pas encore été exécutée
  faute de démon Docker démarré. Le déploiement cluster est documenté mais non
  exécuté.}
\end{frame}

\section{Limites, conclusion et perspectives}

% ================================================================ 18
\begin{frame}{Limites du travail}
  \vspace{2mm}
  \small
  \begin{tabular}{@{}p{0.34\textwidth} p{0.58\textwidth}@{}}
    \toprule
    \textbf{Limite} & \textbf{Précision} \\
    \midrule
    Données synthétiques &
      Générées par le script fourni. Les conclusions démontrent une méthode. \\[1.5mm]
    Mesure exploratoire &
      Un passage par mode, sur un poste partagé, avec 37 \% de variabilité. \\[1.5mm]
    Cluster documenté, non exécuté &
      Commandes HDFS, YARN et \texttt{cluster.conf} figurent au README. \\[1.5mm]
    Détection pédagogique &
      Quatre règles seuillées ne constituent pas un modèle de fraude validé. \\[1.5mm]
    Conteneur non construit &
      Le \texttt{Dockerfile} s'appuie sur les images officielles. \\
    \bottomrule
  \end{tabular}

  \note{Collectif, 60 secondes. Cette diapositive désamorce la moitié des
  questions du jury. Annoncer ces limites nous mêmes vaut mieux que de les
  subir.}
\end{frame}

% ================================================================ 19
\begin{frame}{Bilan et perspectives}
  \vspace{2mm}
  \begin{block}{Bilan}
    Les 138\,047 lignes lues se reconstituent exactement. Chaque rejet porte
    son motif. Le même code passe du poste au cluster sans être réécrit.
  \end{block}

  \vspace{3mm}
  \textbf{Perspectives}
  \vspace{1mm}
  \begin{columns}[t,onlytextwidth]
    \column{0.47\textwidth}
      \footnotesize
      \begin{itemize}\itemsep2mm
        \item Répéter le benchmark en alternant l'ordre des modes
        \item Construire et éprouver l'image du conteneur
      \end{itemize}
    \column{0.47\textwidth}
      \footnotesize
      \begin{itemize}\itemsep2mm
        \item Exécuter le pipeline sur un cluster réel
        \item Mesurer la couverture de lignes des tests
      \end{itemize}
  \end{columns}

  \note{Collectif, 90 secondes. Terminer sur la dernière phrase, c'est le
  message que le jury doit retenir. Enchaîner ensuite directement sur les
  questions. Chaque membre peut être interrogé sur le code d'un autre.}
\end{frame}

% ================================================================ 20
\sanspied
\begin{frame}[plain]
  \vspace{2.6cm}
  \begin{center}
    {\LARGE Merci de votre attention}
  \end{center}

  \note{Réviser en priorité : la différence entre les deux notions de rejet,
  le choix de rangeBetween sur des secondes, et l'exclusion des horizons non
  observables dans le calcul de la rétention.}
\end{frame}

\end{document}
